% 单个ReLU单元 f(x)=a*relu(b*x) 的重缩放：(a,b) 与 (a/c, c b)（c>0）表示同一函数。
% 左：参数平面 x=a, y=b，曲线 ab=1（同一函数 x->relu(x)）；坐标 1 单位 = 14 mm。
% 三点 P1=(2,1/2)、P2=(1,1)、P3=(1/2,2)，a^2+b^2 分别为 4.25、2、4.25；
% 细圆弧 a^2+b^2=2 与曲线相切于 P2（ab 固定时 a^2+b^2>=2ab，等号当且仅当 a=b）。
% 右：三组参数对应的同一函数图像 y=relu(x)。
\begingroup
% 本图使用 specs/figures.md 的低饱和色值；以 lt8 前缀局部定义，避免覆盖共享调色板。
\definecolor{lt8Blue}{HTML}{7188A3}\definecolor{lt8BlueFill}{HTML}{E5EAF0}
\definecolor{lt8Green}{HTML}{829C8B}\definecolor{lt8GreenFill}{HTML}{E8EEE9}
\definecolor{lt8Yellow}{HTML}{B4A16F}\definecolor{lt8YellowFill}{HTML}{F1ECDD}
\definecolor{lt8Red}{HTML}{A7787D}\definecolor{lt8RedFill}{HTML}{F0E5E6}
\definecolor{lt8Ink}{HTML}{3B4252}\definecolor{lt8Muted}{HTML}{6B7280}\definecolor{lt8Stroke}{HTML}{4C4D4F}
\begin{tikzpicture}[figure picture,x=1mm,y=1mm,text=lt8Ink]
 \begin{scope}
  \draw[-{Stealth[length=1.6mm]},lt8Stroke,line width=.6pt] (0,0)--(44,0) node[right] {$a$};
  \draw[-{Stealth[length=1.6mm]},lt8Stroke,line width=.6pt] (0,0)--(0,44) node[above] {$b$};
  \node[below left] at (0,0) {$0$};
  \foreach \v/\l in {14/1,28/2}{\draw[lt8Stroke] (\v,0)--(\v,-1) node[below] {$\l$};
   \draw[lt8Stroke] (0,\v)--(-1,\v) node[left] {$\l$};}
  \draw[lt8Muted,line width=.6pt,dashed] (19.80,0) arc[start angle=0,end angle=90,radius=19.80];
  \draw[lt8Blue,line width=1.1pt,domain=0.33:2.85,samples=80,smooth] plot ({14*\x},{14/\x});
  \node[text=lt8Blue,anchor=west] at (36,9) {$ab=1$};
  \fill[lt8Red] (28,7) circle (1);  \node[anchor=south west] at (28.5,7.5) {$(2,\tfrac12)$};
  \fill[lt8Green] (14,14) circle (1.2); \node[anchor=south west] at (14.8,14.8) {$(1,1)$};
  \fill[lt8Yellow] (7,28) circle (1); \node[anchor=south west] at (7.5,28.5) {$(\tfrac12,2)$};
  \node[figure note,anchor=north] at (21,-8) {同一函数的参数集合是一条曲线};
 \end{scope}
 \begin{scope}[xshift=78mm]
  \draw[-{Stealth[length=1.6mm]},lt8Stroke,line width=.6pt] (-12,0)--(22,0) node[right] {$x$};
  \draw[-{Stealth[length=1.6mm]},lt8Stroke,line width=.6pt] (0,-2)--(0,30) node[above] {$f(x)$};
  \draw[lt8Blue,line width=1.1pt] (-11,0)--(0,0)--(20,20);
  \node[text=lt8Blue,anchor=east] at (-2,10) {$x\mapsto\max\{0,x\}$};
  \node[anchor=north west,align=left] at (3,32) {三组参数输出相同};
  \node[figure note,anchor=north] at (2,-8) {参数平方和为 $4.25$、$2$、$4.25$};
 \end{scope}
\end{tikzpicture}%
\endgroup%
